
\documentclass{article}
\usepackage{amsmath}
\usepackage{amsfonts}

\title{}
\author{}
\date{}

\begin{document}

\maketitle


\subsection{A Simple Model of Bracelets, Calendars, and Distance}
We consider a simple Probit model where deworming behavior is driven by private material incentives - either a Bracelet or a Calendar - and the Distance an individual must walk to the deworming treatment location. The underlying latent variable for the decision to deworm is modeled as:
\[
\begin{aligned}
\text{Deworming}^* = & \ \beta_0 + \beta_1 \cdot \text{Bracelet} + \beta_2 \cdot \text{Calendar} \\
& + \beta_3 \cdot \text{Distance} + \beta_4 \cdot (\text{Bracelet} \times \text{Distance}) \\
& + \beta_5 \cdot (\text{Calendar} \times \text{Distance}) + \epsilon
\end{aligned}
\]
where $\epsilon$ is normally distributed with mean zero. The observable deworming decision is given by:
\[
\text{Deworming} = \mathbb{I}\{\text{Deworming}^* > 0\}
\]
and
\[
\text{Pr}(\text{Deworming}) = \Phi(\text{Deworming}^*)
\]
where $\Phi$ represents the cumulative distribution function (CDF) of the standard normal distribution.

\subsubsection*{Marginal Rate of Substitution (MRS)}

\paragraph{MRS between Bracelet and Distance:}

The MRS between Bracelet and Distance is given by the ratio of the marginal utility of a Bracelet to the marginal utility of Distance:

\[
MRS_{\text{Bracelet, Distance}} = -\frac{\frac{\partial \text{Deworming}^*}{\partial \text{Bracelet}}}{\frac{\partial \text{Deworming}^*}{\partial \text{Distance}}}
\]

Taking the partial derivatives, we have:

\[
\frac{\partial \text{Deworming}^*}{\partial \text{Bracelet}} = \beta_1 + \beta_4 \cdot \text{Distance}
\]
\[
\frac{\partial \text{Deworming}^*}{\partial \text{Distance}} = \beta_3 + \beta_4 \cdot \text{Bracelet} + \beta_5 \cdot \text{Calendar}
\]

Thus, the MRS between Bracelet and Distance is:

\[
MRS_{\text{Bracelet, Distance}} = -\frac{\beta_1 + \beta_4 \cdot \text{Distance}}{\beta_3 + \beta_4 \cdot \text{Bracelet} + \beta_5 \cdot \text{Calendar}}
\]

\paragraph{MRS between Calendar and Distance:}

The MRS between Calendar and Distance is:

\[
MRS_{\text{Calendar, Distance}} = -\frac{\frac{\partial \text{Deworming}^*}{\partial \text{Calendar}}}{\frac{\partial \text{Deworming}^*}{\partial \text{Distance}}}
\]

The relevant partial derivative is:

\[
\frac{\partial \text{Deworming}^*}{\partial \text{Calendar}} = \beta_2 + \beta_5 \cdot \text{Distance}
\]

Hence, the MRS between Calendar and Distance becomes:

\[
MRS_{\text{Calendar, Distance}} = -\frac{\beta_2 + \beta_5 \cdot \text{Distance}}{\beta_3 + \beta_4 \cdot \text{Bracelet} + \beta_5 \cdot \text{Calendar}}
\]

\paragraph{MRS between Bracelet and Calendar:}

The MRS between Bracelet and Calendar is derived similarly:

\[
MRS_{\text{Bracelet, Calendar}} = -\frac{\frac{\partial \text{Deworming}^*}{\partial \text{Bracelet}}}{\frac{\partial \text{Deworming}^*}{\partial \text{Calendar}}}
\]

This simplifies to:

\[
MRS_{\text{Bracelet, Calendar}} = -\frac{\beta_1 + \beta_4 \cdot \text{Distance}}{\beta_2 + \beta_5 \cdot \text{Distance}}
\]

\subsubsection*{Predictions}

Consider the case where the MRS for Bracelet varies with Distance, but the MRS for Calendar remains fixed. Specifically:

\begin{itemize}
    \item The MRS between Bracelet and Distance varies with Distance due to $\beta_4 \neq 0$, indicating that the marginal utility of a Bracelet changes with Distance.
    \item The MRS between Calendar and Distance is constant when $\beta_5 = 0$, meaning that the marginal utility of a Calendar does not vary with Distance.
\end{itemize}

\begin{prediction}
Therefore, the MRS between Bracelet and Calendar is:

\[
MRS_{\text{Bracelet, Calendar}} = -\frac{\beta_1 + \beta_4 \cdot \text{Distance}}{\beta_2}
\]

Since the MRS for Bracelet changes over Distance while the MRS for Calendar is fixed, 
the MRS between Bracelet and Calendar also varies with Distance due to the interaction 
term $\beta_4 \cdot \text{Distance}$.
\end{prediction}

However, the WTP data finds individuals value the private incentive from the Bracelets 
and Calendars identically 
at all distances. In a simple model composed of only private material incentives and distance 
it is impossible for the two private incentives to have differential responses to distance but the 
same MRS across all distances. This suggests that patterns of substitution between the two incentives 
and distance alone are not enough to explain the observed patterns in the data - some other factor,
such as social image concerns, can generate identical marginal utilities for the 
private components of the incentives whilst allowing the overall marginal rate of 
substitution to vary, including both private and signalling utilities.

\end{document}
